\Article{17}
Wie viel Vertrauen dieser Bestimmung beizumessen ist, muss anhand des mittleren Fehlers beurteilt werden, der entweder bei der Bestimmung selbst oder bei ihrem Quadrat zu befürchten ist; im letzteren Fall ist dieser mittlere Fehler die Quadratwurzel
\SourceLeaf{081}
aus dem Mittelwert des Ausdrucks
\[
\left(\frac S\sigma-\mu\mu\right)^2
\]
Seine Entwicklung ergibt sich aus einer ähnlichen Schlussweise wie derjenigen, die in der vorhergehenden Abhandlung, Art. 39 ff., dargelegt wurde. Diese Überlegungen lassen wir hier der Kürze halber aus und geben nur die Formel selbst an. Der bei der Bestimmung des Quadrats \(\mu\mu\) zu befürchtende mittlere Fehler wird nämlich ausgedrückt durch
\[
\BareRoot{\left(\frac{2\mu^4}{\sigma}+\frac{\nu^4-3\mu^4}{\sigma\sigma}\cdot N\right)}
\]
wobei \(\nu^4\) den Mittelwert der vierten Potenzen der Fehler bezeichnet. deren Gewicht \(=1.\) ist, und \(N\) die Summe
\[
\begin{aligned}
&(a\alpha+b\beta+c\gamma+\text{etc.})^2+(a'\alpha'+b'\beta'+c'\gamma'+\text{etc.})^2\\
&\qquad+(a''\alpha''+b''\beta''+c''\gamma''+\text{etc.})^2\text{etc.}
\end{aligned}
\]
Diese Summe lässt sich im Allgemeinen nicht auf eine einfachere Form bringen; doch lässt sich, ähnlich wie in Art. 40 der vorhergehenden Abhandlung gezeigt wurde, ihr Wert stets zwischen \(\pi\) und \(\frac{\sigma\sigma}{\pi}\) einschließen. Unter der Hypothese, auf der die Theorie der kleinsten Quadrate von Anfang an beruhte, fällt der diese Summe enthaltende Term weg, da \(\nu^4=3\mu^4\) gilt; und dem durch die Formel \(\BareRoot{\frac S\sigma}\) bestimmten mittleren Fehler ist dieselbe Genauigkeit beizumessen, als wäre dieser mittlere Fehler aus \(\sigma\) genau bekannten Fehlern nach Art. 15 und 16 der vorhergehenden Abhandlung hergeleitet worden.

\Article{18}
Für die Ausgleichung der Beobachtungen sind, wie wir oben gesehen haben, zwei Dinge erforderlich: Erstens müssen die Korrelate der Bedingungsgleichungen, das heißt die Zahlen \(A,B,C\) etc., die den Gleichungen (12) genügen, abgeleitet werden; zweitens müssen diese Zahlen in die Gleichungen (10) eingesetzt werden. Eine auf diese Weise hervorgehende Ausgleichung kann als \textit{vollkommen} oder \textit{vollständig} bezeichnet werden, um sie von einer \textit{unvollkommenen} oder \textit{unvollständigen} Ausgleichung zu unterscheiden. Mit Letzterer bezeichnen wir Ausgleichungen, die zwar aus denselben Gleichungen (10) hervorgehen, dabei aber von eingesetzten Werten der Größen \(A,B,C\) etc. ausgehen, die den Gleichungen (12) nicht genügen, also nur einigen von ihnen oder gar keiner. Änderungen der Beobachtungen, die sich nicht durch die Formeln (10) erfassen lassen, sind sowohl von der vorliegenden Untersuchung als auch von der Bezeichnung “Ausgleichungen” ausgeschlossen. Da, soweit die Gleichungen (10) gelten, die Gleichungen (13) den Gleichungen (12) völlig gleichwertig sind, lässt sich der Unterschied auch so ausdrücken: \SourceLeaf{082}
Vollständig ausgeglichene Beobachtungen genügen allen Bedingungsgleichungen \(X=0,\allowbreak Y=0,\allowbreak Z=0\) etc.; unvollständig ausgeglichene Beobachtungen genügen entweder keiner oder zumindest nicht allen. Eine Ausgleichung also, die allen Bedingungsgleichungen genügt, ist notwendig vollständig.

\Article{19}
Da aus dem Begriff der Ausgleichung selbst folgt, dass die Kombination zweier Ausgleichungen wiederum eine Ausgleichung bildet, ist leicht einzusehen, dass es keinen Unterschied macht, ob die Regeln zur Bestimmung einer vollkommenen Ausgleichung unmittelbar auf die ursprünglichen Beobachtungen angewandt werden oder auf Beobachtungen, die bereits unvollständig ausgeglichen worden sind.

Es mögen nämlich \(-\theta,-\theta',-\theta''\) etc. ein System unvollständiger Ausgleichungen bilden, das aus den Formeln (I) hervorgegangen ist:
\[
\begin{aligned}
\theta p&=A^0a+B^0b+C^0c+\text{etc.}\\
\theta'p'&=A^0a'+B^0b'+C^0c'+\text{etc.}\\
\theta''p''&=A^0a''+B^0b''+C^0c''+\text{etc.}\\
&\text{etc.}
\end{aligned}
\]
Da vorausgesetzt wird, dass die durch diese Ausgleichungen veränderten Beobachtungen nicht allen Bedingungsgleichungen genügen, seien \(\mathfrak A^*,\mathfrak B^*,\mathfrak C^*\) etc. die Werte, die \(X,Y,Z\) etc. beim Einsetzen jener Beobachtungen annehmen. Zu suchen sind Zahlen \(A^*,B^*,C^*\) etc., die den Gleichungen (II) genügen:
\[
\begin{aligned}
\mathfrak A^*&=A^*[aa]+B^*[ab]+C^*[ac]+\text{etc.}\\
\mathfrak B^*&=A^*[ab]+B^*[bb]+C^*[bc]+\text{etc.}\\
\mathfrak C^*&=A^*[ac]+B^*[bc]+C^*[cc]+\text{etc.}
\end{aligned}
\]
etc. Ist dies geschehen, so wird eine vollständige\textperiodcentered\ Ausgleichung der so veränderten Beobachtungen durch die neuen Änderungen \(-\kappa,-\kappa',-\kappa''\) etc. bewirkt, wobei \(\kappa,\kappa',\kappa''\) etc. nach den Formeln (III) zu berechnen sind:
\[
\begin{aligned}
\kappa p&=A^*a+B^*b+C^*c+\text{etc.}\\
\kappa'p'&=A^*a'+B^*b'+C^*c'+\text{etc.}\\
\kappa''p''&=A^*a''+B^*b''+C^*c''+\text{etc.}
\end{aligned}
\]
Untersuchen wir nun, wie diese Korrekturen mit der vollständigen Ausgleichung der ursprünglichen\SourceLeaf{083}
Beobachtungen zusammenhängen. Zunächst ist offenbar
\[
\begin{aligned}
\mathfrak A^*&=\mathfrak A-a\theta-a'\theta'-a''\theta''-\text{etc.}\\
\mathfrak B^*&=\mathfrak B-b\theta-b'\theta'-b''\theta''-\text{etc.}\\
\mathfrak C^*&=\mathfrak C-c\theta-c'\theta'-c''\theta''-\text{etc.}
\end{aligned}
\]
etc. Setzen wir in diesen Gleichungen für \(\theta,\theta',\theta''\) etc. die Werte aus (I) sowie für \(\mathfrak A^*,\mathfrak B^*,\mathfrak C^*\) etc. die Werte aus (II) ein, so erhalten wir
\[
\begin{aligned}
\mathfrak A&=(A^0+A^*)[aa]+(B^0+B^*)[ab]+(C^0+C^*)[ac]+\text{etc.}\\
\mathfrak B&=(A^0+A^*)[ab]+(B^0+B^*)[bb]+(C^0+C^*)[bc]+\text{etc.}\\
\mathfrak C&=(A^0+A^*)[ac]+(B^0+B^*)[bc]+(C^0+C^*)[cc]+\text{etc.}
\end{aligned}
\]
etc.; daraus folgt, dass die Korrelate der Bedingungsgleichungen, die den Gleichungen (12) genügen, gegeben sind durch
\[
A=A^0+A^*,\quad B=B^0+B^*,\quad C=C^0+C^*\ \text{etc.}
\]
Daraus wiederum folgt aus den Gleichungen (10), I und III,
\[
\varepsilon=\theta+\kappa,\quad\varepsilon'=\theta'+\kappa'\quad\varepsilon''=\theta''+\kappa''\ \text{etc.}
\]
das heißt: Die vollständige Ausgleichung der Beobachtungen ist dieselbe, ob sie unmittelbar berechnet wird oder mittelbar, ausgehend von einer unvollständigen Ausgleichung.